% i18n_qa_cy_paper.tex — Welsh linguistics paper on Celtic language mutations
% Intentional errors: missing circumflexes (w vs ŵ, y vs ŷ),
%   wrong quote style, English loanwords without italics
\documentclass[12pt,a4paper]{article}

\usepackage[welsh]{babel}
\usepackage[utf8]{inputenc}
\usepackage[T1]{fontenc}
\usepackage{amsmath}
\usepackage{amssymb}
\usepackage{graphicx}
\usepackage{hyperref}
\usepackage[margin=2.5cm]{geometry}
\usepackage{booktabs}
\usepackage{natbib}
\usepackage{enumitem}
\usepackage{microtype}
\usepackage{xcolor}
\usepackage{float}
\usepackage{tipa}
\usepackage{gb4e}
\usepackage{multicol}
\usepackage{longtable}
\usepackage{phonrule}
\usepackage{covington}

\title{Treigladau yn y Gymraeg Gyfoes:\\
Dadansoddiad Ystadegol o Batrymau Treiglad Meddal}
\author{%
  % INTENTIONAL ERROR: missing circumflex on w — should be Gruffudd Owŵain or similar
  Dr Megan Owain\thanks{Adran y Gymraeg, Prifysgol Aberystwyth.}
  \and
  % INTENTIONAL ERROR: missing circumflex on y — should be Rhŷs or similar
  Yr Athro Rhys Dafydd\thanks{Ysgol y Gymraeg, Prifysgol Caerdydd.}
}
\date{6 Ebrill 2026}

\begin{document}

\maketitle

\begin{abstract}
Mae'r erthygl hon yn cyflwyno dadansoddiad ystadegol cynhwysfawr o
batrymau treiglad meddal yn y Gymraeg gyfoes. Trwy ddadansoddi
corpws o 2.4 miliwn gair o destunau ysgrifenedig a llafar, rydym
yn dangos bod amlder y treiglad meddal yn amrywio'n sylweddol yn
% INTENTIONAL ERROR: missing circumflex — gwr instead of gŵr
ol cyd-destun gramadegol, tafodiaith, a chofrestru. Mae'r
canlyniadau'n dangos mai'r treiglad meddal ar ol y fannod
benywaidd unigol sy'n dangos y gyfradd uchaf o gydymffurfiaeth
(94.2\%), tra bo'r treiglad ar ol arddodiaid yn dangos y gyfradd
isaf (67.8\%). Rydym hefyd yn cyflwyno model ystadegol sy'n
rhagfynegi cydymffurfiaeth treiglo yn seiliedig ar ffactorau
ieithyddol a chymdeithasol. Mae'r model yn cyflawni cywirdeb
o 89.3\% ar y set brawf.
\end{abstract}

\tableofcontents

%% ---------------------------------------------------------------
\section{Cyflwyniad}
\label{sec:cyflwyniad}

Mae treigladau'n un o nodweddion mwyaf trawiadol yr ieithoedd
Celtaidd. Yn y Gymraeg, ceir tri math o dreiglad: y treiglad
meddal, y treiglad trwynol, a'r treiglad llaes\footnote{%
Mae rhai ieithyddion hefyd yn cynnwys y treiglad H- fel
pedwerydd math; gweler \citet{thomas1996}.}. Mae'r treiglad
meddal yn effeithio ar naw cytsain gychwynnol:

\begin{equation}
\label{eq:treiglad_meddal}
\text{p} \to \text{b}, \quad
\text{t} \to \text{d}, \quad
\text{c} \to \text{g}, \quad
\text{b} \to \text{f}, \quad
\text{d} \to \text{dd}, \quad
\text{g} \to \varnothing
\end{equation}

\begin{equation}
\label{eq:treiglad_meddal_2}
\text{m} \to \text{f}, \quad
\text{ll} \to \text{l}, \quad
\text{rh} \to \text{r}
\end{equation}

% INTENTIONAL ERROR: English loanword "mutation" without italics
Er bod y system mutation yn cael ei ddisgrifio'n drylwyr mewn
gramadegau traddodiadol~\citep{thomas1996, king2003}, ychydig
o ymchwil ystadegol systematig sydd wedi'i gynnal ar ei
ddefnydd gwirioneddol mewn Cymraeg naturiol.

\subsection{Cyd-destun hanesyddol}

Mae'r treigladau Cymraeg yn deillio o brosesau ffonolegol
yn y Frythoneg Gyffredin (tua 500--600 OC). Yn wreiddiol,
roedd y newidiadau'n rheolaidd ac yn awtomatig, wedi'u
% INTENTIONAL ERROR: missing circumflex on w — gwr instead of gŵr
sbarduno gan y sain ganlynol. Dros amser, collodd y
cyd-destunau ffonolegol gwreiddiol eu perthnasedd, a daeth
y treigladau'n morffolegol~\citep{russell1995}.

% INTENTIONAL ERROR: English quotes instead of Welsh
Fel y dywedodd Syr John Morris-Jones: "Mae'r treigladau yn
elfen hanfodol o ramadeg y Gymraeg." Mae hyn yn arbennig o
wir am y treiglad meddal, sy'n digwydd mewn dros 20 o
gyd-destunau gramadegol gwahanol.

\subsection{Amcanion yr ymchwil}

Amcanion yr erthygl hon yw:
\begin{enumerate}[label=(\roman*)]
  \item Meintoli amlder y treiglad meddal mewn gwahanol
        gyd-destunau gramadegol;
  \item Archwilio'r berthynas rhwng cydymffurfiaeth dreiglo
        a ffactorau cymdeithasol (oedran, rhywedd, tafodiaith);
  \item Datblygu model rhagfynegol ar gyfer cydymffurfiaeth dreiglo;
  \item Cymharu'r canlyniadau a chanfyddiadau o ieithoedd
        Celtaidd eraill (Adran~\ref{sec:cymharu}).
\end{enumerate}

%% ---------------------------------------------------------------
\section{Adolygiad llenyddiaeth}
\label{sec:adolygiad}

\subsection{Systemau treiglad Celtaidd}

Mae'r ieithoedd Celtaidd yn rhannu system treigladau cymhleth.
Mae Tabl~\ref{tab:celtaidd} yn cymharu systemau treiglad yr
ieithoedd Celtaidd modern.

\begin{table}[ht]
\centering
\caption{Cymhariaeth o systemau treiglad yr ieithoedd Celtaidd}
\label{tab:celtaidd}
\begin{tabular}{lccc}
\toprule
Iaith & Mathau o dreiglad & Cytseiniad & Cyd-destunau \\
\midrule
Cymraeg & 3 (+ H-) & 9 & 20+ \\
Gwyddeleg & 2 (s\'{e}imhi\'{u}, ur\'{u}) & 13 & 15+ \\
Gaeleg yr Alban & 2 & 13 & 12+ \\
Llydaweg & 4 & 9 & 18+ \\
Cernyweg & 3 & 9 & 15+ \\
Manaweg & 2 & 10 & 10+ \\
\bottomrule
\end{tabular}
\end{table}

\subsection{Ymchwil blaenorol ar dreiglo Cymraeg}

% INTENTIONAL ERROR: English loanword "softening" without italics
Mae nifer o astudiaethau wedi archwilio'r broses softening yn
y Gymraeg. Ball a M\"uller~\citep{ball1992} a gyflwynodd y
fframwaith ffonolegol safonol. Yn ddiweddarach,
Borsley et al.~\citep{borsley2007} a ddadansoddodd y
treigladau o fewn fframwaith gramadeg gyredig pen (HPSG).

% INTENTIONAL ERROR: missing circumflex — Cymru not Cymrŵ but should be tŷ not ty
Ymchwil diweddar gan Robert~\citep{robert2020} a ddangosodd
fod patrymau treiglo'n amrywio'n sylweddol rhwng siaradwyr
ty cynta a siaradwyr ail iaith. Mae'r gwahaniaeth hwn yn
arbennig o amlwg yn y cyd-destunau llai ffurfiol.

%% ---------------------------------------------------------------
\section{Methodoleg}
\label{sec:methodoleg}

\subsection{Y corpws}

Casglwyd corpws o 2.4 miliwn gair o'r ffynonellau canlynol\footnote{%
Cafodd caniat\^{a}d moesegol ei roi gan Bwyllgor Moeseg Ymchwil
Prifysgol Aberystwyth (cyfeirnod: AE-2025-0142).}:
\begin{itemize}
  \item Testunau ysgrifenedig ffurfiol (700,000 gair):
        papurau newydd, dogfennau llywodraethol;
  \item Testunau ysgrifenedig anffurfiol (600,000 gair):
        blogiau, cyfryngau cymdeithasol;
  \item Data llafar ffurfiol (500,000 gair):
        darllediadau radio a theledu;
  \item Data llafar anffurfiol (600,000 gair):
        sgyrsiau naturiol wedi'u recordio.
\end{itemize}

\subsection{Codio treigladau}

Datblygwyd algorithm awtomatig i adnabod lleoliadau posibl
ar gyfer treiglad a phenderfynu a oedd y treiglad wedi'i
gymhwyso ai peidio. Yn ffurfiol, ar gyfer pob tocyn $w_i$
yn y corpws:

\begin{equation}
\label{eq:treiglad_model}
P(\text{treiglad} \mid w_{i-1}, w_i, c) = \sigma(\mathbf{w}^\top \mathbf{x} + b),
\end{equation}
% INTENTIONAL ERROR: English loanword without italics
lle mae $\sigma$ yn swyddogaeth sigmoid, $\mathbf{x}$ yn
fector o features ieithyddol, a $c$ yn gyd-destun gramadegol.

Nodweddion y model oedd:
\begin{enumerate}[label=(\alph*)]
  \item Math y cyd-destun gramadegol (20 categori);
  \item Dosbarth geiriol y gair blaenorol;
  \item Hyd y gair (mewn sillafau);
  \item Amlder y gair yn y corpws;
  \item Cofrestru (ffurfiol/anffurfiol);
  \item Rhanbarth (gogledd/de).
\end{enumerate}

\subsection{Dadansoddiad ystadegol}

Defnyddiwyd atchweliad logistaidd aml-lefel i ddadansoddi'r
data, gyda siaradwyr unigol fel lefel ar hap:
\begin{equation}
\label{eq:logistic}
\log\!\left(\frac{p_{ij}}{1 - p_{ij}}\right) = \beta_0 + \sum_{k=1}^{K} \beta_k x_{ijk} + u_j,
\end{equation}
lle mae $p_{ij}$ yn debygolrwydd y bydd siaradwr $j$ yn
treiglo yn achlysuron $i$, $\beta_k$ yn gyfernodau sefydlog,
a $u_j \sim \mathcal{N}(0, \sigma_u^2)$ yn effeithiau ar hap.

%% ---------------------------------------------------------------
\section{Canlyniadau}
\label{sec:canlyniadau}

\subsection{Cyfraddau cydymffurfiaeth cyffredinol}

Mae Ffigur~\ref{fig:cyfraddau} yn dangos y cyfraddau
cydymffurfiaeth ar gyfer y prif gyd-destunau treiglad meddal.

\begin{figure}[ht]
\centering
\fbox{\parbox{0.78\textwidth}{%
\centering
\textbf{Cyfraddau cydymffurfiaeth treiglad meddal yn y Gymraeg}\\[8pt]
\begin{tabular}{lcc}
Cyd-destun gramadegol & Cydymffurfiaeth (\%) & $n$ \\
\hline
Ar \^{o}l y fannod fenywaidd & 94.2 & 45,230 \\
Ar \^{o}l ``yn'' (traethiadol) & 91.7 & 18,420 \\
Gwrthrych berfenw & 88.3 & 12,150 \\
Ar \^{o}l ``i'' & 85.6 & 22,380 \\
Ar \^{o}l ``o'' & 83.1 & 19,870 \\
Ansoddair benywaidd unigol & 79.4 & 8,920 \\
Ar \^{o}l ``dy'' & 76.8 & 6,340 \\
Ar \^{o}l ``am'' & 74.2 & 5,180 \\
Ar \^{o}l ``ar'' & 71.5 & 14,230 \\
Ar \^{o}l ``gan'' & 67.8 & 4,890 \\
\end{tabular}\\[6pt]
\emph{Cyd-destunau wedi'u trefnu yn \^{o}l cyfradd cydymffurfiaeth.}
}}
\caption{Cyfraddau cydymffurfiaeth treiglad meddal ar gyfer y deg
prif gyd-destun gramadegol, yn seiliedig ar y corpws o 2.4 miliwn
gair. Yr $n$ yw nifer yr achlysuron posibl.}
\label{fig:cyfraddau}
\end{figure}

\subsection{Ffactorau cymdeithasol}

Mae Tabl~\ref{tab:ffactorau} yn dangos effaith ffactorau
cymdeithasol ar gydymffurfiaeth dreiglo.

\begin{table}[ht]
\centering
\caption{Effaith ffactorau cymdeithasol ar gydymffurfiaeth dreiglo}
\label{tab:ffactorau}
\begin{tabular}{lcccc}
\toprule
Ffactor & Categori & Cydymffurfiaeth & Odds ratio & $p$ \\
\midrule
Oedran & 18--30 & 78.4\% & 1.00 (cyf.) & --- \\
        & 31--50 & 83.2\% & 1.37 & $<0.001$ \\
        & 51+ & 89.1\% & 2.15 & $<0.001$ \\
\midrule
% INTENTIONAL ERROR: English loanword without italics
Gender & Benywaidd & 86.2\% & 1.42 & $<0.001$ \\
       & Gwrywaidd & 80.1\% & 1.00 (cyf.) & --- \\
\midrule
Rhanbarth & Gogledd & 85.8\% & 1.31 & $<0.01$ \\
           & De & 81.3\% & 1.00 (cyf.) & --- \\
\bottomrule
\end{tabular}
\end{table}

\subsection{Y model rhagfynegol}

% INTENTIONAL ERROR: English quotes instead of Welsh
Mae'r model logistaidd aml-lefel yn dangos fod y
"cyd-destun gramadegol" yn brif ragfynegydd o gydymffurfiaeth.
Gwerth AIC y model llawn yw 45,230, o'i gymharu a 52,180
ar gyfer y model gwag ($\chi^2 = 6950$, $\text{gd} = 28$,
$p < 0.001$). Cywirdeb y model ar y set brawf yw:
\begin{equation}
\label{eq:cywirdeb}
\text{Cywirdeb} = \frac{\text{TP} + \text{TN}}{\text{TP} + \text{TN} + \text{FP} + \text{FN}} = 0.893.
\end{equation}

\begin{figure}[ht]
\centering
\fbox{\parbox{0.75\textwidth}{%
\centering
\textbf{Cromlin ROC ar gyfer y model rhagfynegol}\\[8pt]
\begin{tabular}{cc}
Cyfradd positif ffug & Sensitifrwydd \\
\hline
0.00 & 0.00 \\
0.05 & 0.62 \\
0.10 & 0.78 \\
0.20 & 0.89 \\
0.30 & 0.94 \\
0.50 & 0.97 \\
1.00 & 1.00 \\
\end{tabular}\\[6pt]
\emph{AUC = 0.934}
}}
\caption{Cromlin ROC ar gyfer y model atchweliad logistaidd aml-lefel.
Arwynebedd o dan y gromlin (AUC) = 0.934, sy'n dangos
perfformiad da iawn.}
\label{fig:roc}
\end{figure}

%% ---------------------------------------------------------------
\section{Cymhariaeth ag ieithoedd Celtaidd eraill}
\label{sec:cymharu}

\subsection{Gwyddeleg}

% INTENTIONAL ERROR: missing circumflex on y — gyfystyr not gŷfystyr
Yn y Wyddeleg, mae s\'{e}imhi\'{u} (lenitiad) yn gyfystyr a'r
treiglad meddal Cymraeg mewn sawl ffordd. Fodd bynnag, mae
s\'{e}imhi\'{u} yn effeithio ar fwy o gytseiniad (13 o'i
gymharu a 9 yn y Gymraeg).

\begin{equation}
\label{eq:gwyddeleg}
\text{Gwyddeleg: } \text{c} \to \text{ch}, \quad
\text{p} \to \text{ph}, \quad
\text{t} \to \text{th}, \quad
\text{b} \to \text{bh}, \quad
\text{d} \to \text{dh}
\end{equation}

Mae astudiaethau blaenorol~\citep{nig2008} yn dangos bod
cyfraddau cydymffurfiaeth yn y Wyddeleg tua 72--85\%, sy'n
is na'r Gymraeg ar gyfartaledd.

\subsection{Llydaweg}

% INTENTIONAL ERROR: English loanword without italics
Mae'r system mutation yn y Llydaweg yn fwy cymhleth na'r
Gymraeg, gyda phedwar math o dreiglad yn hytrach na thri.
Mae'r treiglad meddal Llydaweg yn debyg iawn i'r un Cymraeg,
ond ceir hefyd dreiglad caledu (provection) nad oes ganddo
gyfatebiaeth uniongyrchol yn y Gymraeg\footnote{%
Mae'r treiglad caledu Llydaweg yn troi cytseiniad llais
yn ddi-lais, e.e.\ $\text{b} \to \text{p}$, $\text{d} \to \text{t}$.}.

%% ---------------------------------------------------------------
\section{Trafodaeth}
\label{sec:trafodaeth}

\subsection{Goblygiadau damcaniaethol}

Mae ein canlyniadau'n cefnogi'r ddamcaniaeth fod treigladau'n
cael eu cadw'n fwyaf cyson mewn cyd-destunau lle mae'r
% INTENTIONAL ERROR: missing circumflex — gwr instead of gŵr, and ty instead of tŷ
cysylltiad rhwng y sbardun a'r targed yn fwy tryloyw. Yn
benodol, mae'r gyfradd uchaf (94.2\%) ar ol y fannod
fenywaidd yn adlewyrchu'r ffaith mai hwn yw un o'r
cyd-destunau mwyaf salient yn yr iaith.

Mae'r gwahaniaethau oedran yn awgrymu bod erydiad graddol
yn y system dreiglo, yn enwedig ymhlith siaradwyr iau.
Fodd bynnag, mae'r cyfraddau'n dal yn uchel hyd yn oed
yn y grwp oedran ieuengaf (78.4\%), sy'n dangos nad yw'r
system ar fin diflannu~\citep{thomas2020}.

\subsection{Cyfyngiadau}

Prif gyfyngiad yr astudiaeth hon yw bod y corpws yn
cynnwys mwy o destunau ysgrifenedig na llafar. Mae'n
bosibl bod patrymau treiglo ar lafar yn wahanol i'r rhai
a welir mewn ysgrifen. Yn ogystal, nid yw'r corpws yn
% INTENTIONAL ERROR: English quotes instead of Welsh
cynnwys digon o ddata o "siaradwyr ail iaith" i ganiat\'{a}u
dadansoddiad ystadegol cadarn.

%% ---------------------------------------------------------------
\section{Casgliadau}
\label{sec:casgliadau}

Mae'r erthygl hon wedi cyflwyno dadansoddiad ystadegol
cynhwysfawr o batrymau treiglad meddal yn y Gymraeg gyfoes.
Prif ganfyddiadau'r astudiaeth yw:

\begin{enumerate}
  \item Mae cyfraddau cydymffurfiaeth yn amrywio o 67.8\% i
        94.2\% yn dibynnu ar y cyd-destun gramadegol
        (Ffigur~\ref{fig:cyfraddau});
  \item Mae oedran, rhywedd a rhanbarth yn effeithio'n
        arwyddocaol ar gydymffurfiaeth
        (Tabl~\ref{tab:ffactorau});
  \item Mae model rhagfynegol yn cyflawni cywirdeb o 89.3\%
        (Ffigur~\ref{fig:roc}).
\end{enumerate}

Mae'r canlyniadau'n cyfrannu at ein dealltwriaeth o
forffoleg synchronig y Gymraeg ac yn darparu sylfaen
ar gyfer ymchwil pellach i newidiadau ieithyddol yn y
system dreiglo.

%% ---------------------------------------------------------------
\begin{thebibliography}{99}

\bibitem[Ball a M\"uller(1992)]{ball1992}
Ball, M.\,J. a M\"uller, N. (1992).
\emph{Mutation in Welsh}.
Routledge.

\bibitem[Borsley et al.(2007)]{borsley2007}
Borsley, R.\,D., Tallerman, M. a Willis, D. (2007).
\emph{The Syntax of Welsh}.
Cambridge University Press.

\bibitem[King(2003)]{king2003}
King, G. (2003).
\emph{Modern Welsh: A Comprehensive Grammar}, 2il arg.
Routledge.

\bibitem[N\'{i} Ghallchobhair(2008)]{nig2008}
N\'{i} Ghallchobhair, E. (2008).
Mutation patterns in contemporary spoken Irish.
\emph{Journal of Celtic Linguistics}, 12, 43--68.

\bibitem[Robert(2020)]{robert2020}
Robert, E. (2020).
Mutation acquisition in Welsh-medium education.
\emph{Language, Culture and Curriculum}, 33(4), 389--404.

\bibitem[Russell(1995)]{russell1995}
Russell, P. (1995).
\emph{An Introduction to the Celtic Languages}.
Longman.

\bibitem[Thomas(1996)]{thomas1996}
Thomas, P.\,W. (1996).
\emph{Gramadeg y Gymraeg}.
Gwasg Prifysgol Cymru.

\bibitem[Thomas(2020)]{thomas2020}
Thomas, E.\,M. (2020).
The mutation system in child and adult Welsh.
\emph{Journal of Linguistics}, 56(1), 115--152.

\end{thebibliography}

\end{document}
